\chapter{Laparoscopy (overview)}
\index{Laparoscopy (overview)}

\section*{Definition and Overview}
Laparoscopy, also known as keyhole surgery or minimally invasive surgery, is a modern surgical technique that allows a surgeon to access the interior of the abdomen and pelvis without having to make large abdominal incisions. The procedure utilizes a laparoscope --- a long, thin tube equipped with a high-intensity light source and a high-resolution digital camera at its tip. The laparoscope is inserted through a small skin incision (typically 5 to 12 mm), transmitting real-time, magnified video images of the internal organs to a monitor in the operating room. To create a working space, the abdominal cavity is inflated (insufflated) with carbon dioxide gas. Laparoscopy is widely used for both diagnostic and therapeutic purposes across various medical specialties, including general surgery, gynaecology, and urology.

\section*{Epidemiology}
Laparoscopy has become the global standard of care for numerous abdominal and pelvic procedures. Annually, tens of millions of laparoscopic operations are performed worldwide. In many countries, it is the primary surgical modality for cholecystectomies, appendicectomies, bariatric surgeries, and the management of gynaecological conditions such as endometriosis and ectopic pregnancies. The adoption of laparoscopy has significantly reduced the average length of hospital stays and post-operative recovery times across all age demographics.

\section*{Aetiology and Pathophysiology}
Laparoscopy is selected as a surgical approach due to its physiological benefits over open laparotomy. The underlying concept is the minimisation of surgical trauma while achieving the same therapeutic outcome. Indications and physiological features include:

\begin{itemize}
    \item \textbf{Diagnostic laparoscopy:} Indicated for the investigation of chronic pelvic pain, staging of abdominal cancers, or evaluating unexplained abdominal masses when non-invasive imaging is inconclusive.
    \item \textbf{Therapeutic laparoscopy:} Indicated for organ resection (e.g.\ laparoscopic cholecystectomy, appendicectomy, colectomy), hernia repairs, and anti-reflux procedures.
    \item \textbf{Reduced inflammatory response:} Minimised tissue handling and smaller incisions lead to a decreased systemic inflammatory response, resulting in less tissue oedema and postoperative pain.
    \item \textbf{Pneumoperitoneum creation:} Insufflation of the abdomen with carbon dioxide gas to a pressure of 12--15 mmHg. This gas is chosen because it is non-flammable, easily absorbed by the body, and eliminated via the lungs, though it can cause temporary diaphragmatic irritation.
\end{itemize}

\section*{Clinical Features}
A patient recovering from a laparoscopic procedure will exhibit clinical features characteristic of minimally invasive access. These features include:

\begin{itemize}
    \item \textbf{Port-site incisions:} Small surgical wounds on the abdomen (usually 1 to 4 incisions), closed with dissolvable sutures, staples, or surgical adhesive.
    \item \textbf{Referred shoulder pain:} A common postoperative feature caused by residual carbon dioxide gas irritating the phrenic nerve, which shares nerve pathways with the shoulder.
    \item \textbf{Mild abdominal distension:} Temporary bloating and fullness resulting from the pneumoperitoneum.
    \item \textbf{Anaesthetic and gas side effects:} Mild nausea, dizziness, and throat discomfort from the endotracheal tube used during general anaesthesia.
\end{itemize}

\section*{Differential Diagnosis}
In the postoperative period, clinicians must differentiate normal recovery symptoms from serious complications. Key considerations include:

\begin{itemize}
    \item \textbf{Referred shoulder pain vs.\ myocardial infarction:} Distinguishing gas-induced phrenic nerve irritation from acute cardiac ischaemia by obtaining an ECG and monitoring for chest tightness or radiation down the arm.
    \item \textbf{Normal postoperative ileus vs.\ bowel injury:} Differentiating temporary bowel sluggishness from mechanical bowel obstruction or peritonitis caused by inadvertent thermal or mechanical bowel perforation.
    \item \textbf{Port-site inflammation vs.\ wound infection:} Differentiating the normal early inflammatory response of healing incisions from surgical site infections characterized by purulent discharge and spreading erythema.
    \item \textbf{Postoperative pain vs.\ internal haemorrhage:} Distinguishing expected incisional pain from severe pain associated with tachycardia, hypotension, and pallor, which suggests active intra-abdominal bleeding.
\end{itemize}

\section*{When to Seek Medical Care}
Patients recovering from laparoscopy at home should monitor their symptoms closely and understand when to seek medical evaluation.

\textbf{Contact your local emergency services for:}
\begin{itemize}
    \item Severe, rapidly worsening abdominal pain that is unresponsive to prescribed analgesia.
    \item Signs of severe internal bleeding, such as sudden dizziness, fainting, extreme pale skin, or rapid heart rate.
    \item Difficulty breathing, shortness of breath, or chest pain.
    \item Inability to tolerate oral fluids accompanied by persistent, uncontrollable vomiting.
    \item High fever (above 38.5 degrees Celsius) with chills or rigors.
\end{itemize}

\section*{Diagnosis and Investigations}
Preparation for laparoscopy requires a comprehensive preoperative assessment to ensure patient safety under general anaesthesia. Diagnostic steps include:

\begin{itemize}
    \item \textbf{Pre-operative blood panels:} Full blood count (FBC) to check for anaemia or baseline infection, electrolytes, urea, creatinine, and coagulation profiles.
    \item \textbf{Pre-operative imaging:} Ultrasound, CT scans, or MRI of the abdomen/pelvis to define the anatomical pathology and plan optimal port placement.
    \item \textbf{Electrocardiogram (ECG):} Recommended for older patients or those with a history of cardiovascular disease to assess cardiac fitness.
    \item \textbf{Pregnancy test:} Conducted routinely in females of childbearing potential to avoid administering general anaesthesia during early pregnancy.
\end{itemize}

\section*{Management}
Laparoscopy is performed in an operating theatre under strict sterile conditions. The management pathway comprises:

\begin{itemize}
    \item \textbf{General anaesthesia:} Essential to ensure complete muscle relaxation, airway protection, and patient comfort during abdominal insufflation.
    \item \textbf{Surgical procedure:} Insertion of a primary trocar (usually at the umbilicus) to establish pneumoperitoneum, followed by insertion of secondary trocars under direct vision to introduce surgical instruments.
    \item \textbf{Postoperative pain management:} Multimodal analgesia utilising paracetamol, NSAIDs, and local anaesthetic infiltration at the port sites during wound closure.
    \item \textbf{Early ambulation:} Encouraging patients to walk within hours of surgery to promote gastrointestinal motility, facilitate carbon dioxide absorption, and prevent venous thromboembolism.
    \item \textbf{Wound care:} Keeping port sites dry for the first 24--48 hours, followed by gentle washing and monitoring of the healing incisions.
\end{itemize}

\section*{Complications}
Although laparoscopy has a lower complication rate than open surgery, serious risks exist and require prompt recognition:

\begin{itemize}
    \item \textbf{Organ injury:} Inadvertent puncture, laceration, or thermal injury to the bowel, bladder, ureters, or major retroperitoneal blood vessels (such as the aorta or vena cava).
    \item \textbf{Gas embolism:} A rare but catastrophic event where carbon dioxide gas enters the venous circulation, causing cardiovascular collapse.
    \item \textbf{Port-site hernia:} Herniation of abdominal contents (e.g.\ omentum or bowel loop) through a trocar incision site, most commonly occurring at ports larger than 10 mm.
    \item \textbf{Subcutaneous emphysema:} Accumulation of carbon dioxide gas in the subcutaneous tissues of the abdominal wall, which typically resolves spontaneously.
    \item \textbf{Conversion to laparotomy:} The necessity to switch to an open surgical incision due to dense adhesions, anatomical obscurity, or intra-operative complications.
\end{itemize}

\section*{Prognosis and Outcomes}
The prognosis for patients undergoing laparoscopic surgery is excellent. Compared to open laparotomy, laparoscopy offers reduced postoperative pain, minimal scarring, a lower risk of incisional hernia and wound infection, and a shorter hospital stay (often as a day-case or overnight stay). Most patients can return to light activities within a few days and resume normal daily activities, including returning to work, within 1 to 2 weeks, depending on the complexity of the underlying procedure.

\section*{Prevention and Self-Care}
To ensure optimal recovery and prevent complications after laparoscopy, patients should:

\begin{itemize}
    \item \textbf{Mobilise frequently:} Walk gently around the house to help disperse residual carbon dioxide gas and lower the risk of deep vein thrombosis.
    \item \textbf{Follow dietary guidelines:} Start with light, easily digestible meals (such as soups and crackers) before returning to a regular diet to minimise nausea and bloating.
    \item \textbf{Avoid heavy lifting:} Restrict lifting objects heavier than 5 kg for the first 2 to 4 weeks to prevent port-site herniation.
    \item \textbf{Maintain wound hygiene:} Clean the incisions daily with mild soap and water once the initial dressing is removed, drying the area thoroughly.
    \item \textbf{Manage bowel habits:} Drink plenty of water and use mild stool softeners if needed to prevent straining due to opioid-induced constipation.
\end{itemize}

\section*{Sources and Further Reading}
The following organisations provide reputable educational materials on laparoscopic surgery:

\begin{itemize}
    \item Healthdirect Australia. \textit{Overview of keyhole surgery, including what to expect and recovery tips.} https://www.healthdirect.gov.au/
    \item NHS. \textit{Comprehensive guide to laparoscopy, explaining the procedure and potential risks.} https://www.nhs.uk/
    \item Mayo Clinic. \textit{Clinical overview of laparoscopy, highlighting diagnostic and therapeutic uses.} https://www.mayoclinic.org/
    \item Royal Australasian College of Surgeons. \textit{Patient resource sheets detailing laparoscopic procedures and surgical safety.} https://www.surgeons.org/
    \item Royal Australian and New Zealand College of Obstetricians and Gynaecologists. \textit{Patient guide to gynaecological laparoscopy.} https://ranzcog.edu.au/
\end{itemize}
